\section{Evaluation}
\label{sec:evaluation}

We evaluate \sysname on four axes: unlearning effectiveness (E6),
on-chain cost on PureChain versus Polygon Amoy (E7), scalability of
round validation to $N\!=\!500$ (E8), and utility recovery through
revocation; two short on-chain demos additionally validate the
liveness and verifiability properties. The full
reproducibility recipe is in supplementary
\texttt{docs/eval-protocol.md}.

\subsection{Experimental setup}
\label{sec:eval-setup}

The primary federation is a 9-client N-BaIoT autoencoder anomaly
detector~\cite{meidan2018nbaiot}, one client per IoT device; the global
model is a symmetric MLP autoencoder
($115\!\to\!75\!\to\!50\!\to\!25\!\to\!10$) trained with FedAvg over 30
rounds (Flower~\cite{beutel2020flower}, seed 42). At round 30 we
simulate a GDPR Art.~17 erasure against client 5 and trigger one of the
three approximate methods or, as the exact reference, full 8-client
retraining from scratch. On-chain experiments run against live
contracts on gas-free PureChain (PoA$^2$) and gas-based Polygon Amoy
(PoS), with $\mathit{VC}$ five ECDSA keys, $t=3$, and $\dmax=1000$
blocks; every run records an \texttt{env.json} fingerprint (commit,
library versions, hardware, seed).

\subsection{Unlearning effectiveness}
\label{sec:eval-unlearning}

Table~\ref{tab:t1} reports each method's three independent MIA scores
against the revoked client's member set; the on-chain
\texttt{submitAttestation} predicate accepts a UA iff all three sit
below $\tauMIA\!=\!0.55$.

\begin{table}[t]
\centering
\caption{Unlearning effectiveness (E6). MIA scores
against the revoked client's training data, evaluated on the
post-unlearning model. The on-chain gate accepts a UA iff all
three are below $\tauMIA = 0.55$.}
\label{tab:t1}
\footnotesize
\setlength{\tabcolsep}{4pt}
\resizebox{\columnwidth}{!}{% Table 1 — N-BaIoT unlearning effectiveness (LiRA α=0.05 fixed)
% generated by fl/experiments/collect_results.py
\begin{tabular}{lcccc}
\toprule
Method & Shadow MIA & LiRA & Label-only & On-chain accept \\
\midrule
No unlearning (baseline) & 0.5455 & 0.0475 & 0.0060 & \checkmark \\
Gradient ascent & 0.4380 & 0.0435 & 0.0030 & \checkmark \\
SISA-Fed & 0.2542 & 0.0795 & 0.0862 & \checkmark \\
FedEraser & 0.5232 & 0.0325 & 0.0050 & \checkmark \\
Naive retrain (exact ref.) & 0.5843 & 0.0445 & 0.0138 & $\times$ \\
\bottomrule
\end{tabular}
}
\end{table}

Three observations stand out. All three approximate methods produce
UAs the chain accepts (shadow $0.438$/$0.254$/$0.523$ for gradient
ascent/SISA-Fed/FedEraser; LiRA and label-only well below $\tauMIA$).
The hero result is that the naive-retrain baseline---a model that
\emph{never saw} client 5's data---yields shadow $0.584>\tauMIA$ and is
\emph{rejected}: shadow picks up client 5's distinctive traffic
distribution, not membership, a known confound~\cite{shokri2017mia}
that LiRA ($0.045$) and label-only ($0.014$) do not exhibit; the gate
is thus faithful to MIA evidence, not the controller's claim
(Section~\ref{sec:discussion}). Gradient ascent's shadow even sits
\emph{below} naive retrain's---anti-fitting moves the model further
from client 5's distribution than never seeing it---tracing a
forgetting--utility frontier (SISA-Fed strong-forgetting; gradient
ascent/FedEraser preserve-utility). All three UAs executed end-to-end
on PureChain, finalising the token to \textsc{Unlearned} in the
attestation transaction (records in supplementary
\texttt{onchain\_attestation.json}).

\subsection{On-chain cost and scalability}
\label{sec:eval-gas}

Table~\ref{tab:t2} reports per-operation gas usage and cost from a
lifecycle ×10 variance run (E7; USD at POL/USD $=0.095347$, CoinGecko,
measurement day).

\begin{table}[t]
\centering
\caption{Per-operation gas and cost (E7), lifecycle ×10
mean. PureChain costs are exactly $\$0$; Amoy costs are non-zero.
Per-DS lifecycle = sum of the four ops.}
\label{tab:t2}
\footnotesize
\setlength{\tabcolsep}{3pt}
\resizebox{\columnwidth}{!}{% Table 2 — Per-operation gas + USD (lifecycle ×10 means)
% generated by fl/experiments/collect_results.py
\begin{tabular}{lrrrr}
\toprule
Operation & PureChain gas & Amoy gas & Amoy MATIC & Amoy USD \\
\midrule
\texttt{updateModelHash} & 30,074 & 32,673 & 0.0010 & \$0.00009 \\
\texttt{grantConsent} & 144,868 & 153,568 & 0.0046 & \$0.00044 \\
\texttt{revokeConsent} & 107,904 & 96,404 & 0.0029 & \$0.00028 \\
\texttt{submitAttestation} (3-of-5) & 102,903 & 137,600 & 0.0041 & \$0.00039 \\
\midrule
\textbf{Per-DS lifecycle} & \textbf{385,750} & \textbf{420,245} & \textbf{0.0126} & \textbf{\$0.00120} \\
\bottomrule
\end{tabular}
}
\end{table}

Per-data-subject lifecycle cost is $\$0$ on PureChain and $\$0.00120$
on Amoy. Gas counts differ modestly between the two chains, but the
dominant factor at scale is gas \emph{price} (zero versus
$\sim\!30$ gwei), not gas count.

$\texttt{validateRound}(N)$ scales exactly linearly in the participant
count. PureChain runs the full $N\!\in\![5,500]$ sweep live
($R^2\!\approx\!1$); Amoy, measured at $N\!\in\!\{5,10\}$ and
extrapolated linearly, projects to $\$0.0176$ per $N\!=\!500$ round.
At a typical 100-round/day cadence that is $\$1.76$/day per deployment
on Amoy versus $\$0$ on PureChain: gas-free infrastructure is a
different cost class, not an optimisation.
Two further PureChain runs confirm the formal properties empirically: a
$\dmax\!=\!2$ timeout fires \texttt{ComplianceViolation} (liveness),
and a $t$-of-5 signature sweep exercises both branches of the
acceptance predicate (verifiability).

\subsection{Utility recovery after revocation}
\label{sec:eval-utility}

Resuming FL for ten rounds after dropping client 5, the post-revoke
model reaches reconstruction error $0.4720$, $6.1\%$ \emph{better} than
the 9-client baseline ($0.5025$)---client 5 was an outlier whose
removal tightens the fit. Gradient ascent ($0.4726$) and FedEraser
($0.4698$) match this at no utility cost, while SISA-Fed pays $+29\%$
($0.6081$ by round 40), consistent with its strongest-forgetting
position and lowest shadow MIA in Table~\ref{tab:t1}.

\subsection{Cross-domain validation (CASAS)}
\label{sec:eval-casas}

On CASAS (5-client BiLSTM), the distribution-shift confound is sharper
still: shadow rejects every method but SISA-Fed
($0.499$/$0.035$/$0.000$), the sole gate-passer, consistent with
N-BaIoT---full cross-domain results appear in the journal version.
